\subsection{良过滤模的希尔伯特–塞缪尔级数}

在本段余下部分，我们使用如下记号：若 \(\mathbf{G} \in \mathbf{Z}((\mathbf{T}))\) 且 \(r \in \mathbf{N}\)，置 \(\mathbf{G}^{(r)} = (1 - \mathbf{T})^{-r}\mathbf{G}\)；特别地，若 \(\mathbf{G} = \sum_{n \in \mathbf{Z}} a_n \mathbf{T}^n\)，则

\[
\mathrm{G} ^ {(1)} = \sum_ {n \in \mathbf {Z}} \left(\sum_ {i \leqslant n} a _ {i}\right) \mathrm{T} ^ {n}.
\]

若 \(\mathbf{G} \geqslant 0\)，则有 \(\mathbf{G}^{(r)} \geqslant 0\) 对每个 \(r \in \mathbf{N}\) 成立。。

设 A 为诺特环，q 为 A 的理想，M 为有限生成 A-模。回忆（III, § 3, no. 1, def. 1），M 上的 q-良过滤是从 Z 到 M 的子模集合的映射 F : n ↦ Fn，满足以下三个条件：

\begin{enumerate}
\def\labelenumi{\alph{enumi})}
\item
  有 \(\mathfrak{q}F_n\subset \mathrm{F}_{n + 1}\subset \mathrm{F}_n\) 对每个 \(n\in \mathbf{Z},\)
\item
  存在 \(n_0 \in \mathbf{Z}\)，使得 \(\mathfrak{q}F_n = \mathrm{F}_{n+1}\) 当 \(n \geqslant n_0\) 时成立；
\item
  存在 \(n_{1} \in \mathbf{Z}\)，使得 \(\mathrm{F}_{n_1} = \mathbf{M}\)。
\end{enumerate}

若 \(n_{0}\) 和 \(n_{1}\) 满足上述条件，则对于所有 \(n\in\mathbf{Z}\)，有

\[
\mathfrak {q} ^ {n - n _ {1}} \mathrm{M} \subset \mathrm{F} _ {n} \subset \mathfrak {q} ^ {n - n _ {0}} \mathrm{M}\tag{7}
\]

（回忆一下，按约定，置 \(\mathfrak{q}^r = \mathbf{A}\)，其中 \(r\leqslant 0\)。）

引理 3.--- 若 F 和 F' 是 M 上的两个 q-良过滤，则存在整数 m，使得 \(F_{n}^{\prime} \subset F_{n-m}\) 对每个 \(n \in Z\) 成立。

事实上，存在 \(n_{2}\)，使得 \(\mathbf{F}_{n}^{\prime}\subset\mathfrak{q}^{n-n_{2}}\mathbf{M}\) 对每个 \(n\) 成立，因而 \(\mathbf{F}_{n}^{\prime}\subset\mathbf{F}_{n-(n_{2}-n_{1})}\) 对每个 \(n\) 成立。

引理 4. --- 设 F 是 M 上的 q-良过滤。若 M/qM 具有有限长度，则对于所有 n ∈ Z，M/F \(_{n+1}\) 和 F \(_{n}\) /F \(_{n+1}\) 都具有有限长度。

用 (7) 中的记号，有 long(M/F \(_{n+1}\) ) ≤ long(M/q \(^{n-n_1+1}\) M)，只需证明对于每个 n，q \(^{n}\) M/q \(^{n+1}\) M 具有有限长度。因此可归结为 q-进过滤的情形。设 (x \(_{1}\) , \ldots, x \(_{r}\) ) 是 A-模 q 的有限生成系，I 是 r 个变量中总次数为 n 的单项式组成的有限集

\(X_{1},...,X_{r}\)。从 \((\mathbf{M}/\mathfrak{q}\mathbf{M})^{\mathbf{I}}\) 到 \(q^{n}M/q^{n+1}M\) 的同态把族 \((u_{m})_{m\in\mathbf{I}}\) 映到元素 \(\sum m(x_{1},...,x_{r})u_{m}\)，并且是满射。由于 M/qM 具有有限长度，\(q^{n}M/q^{n+1}M\) 也具有有限长度。

现在假设 M/qM 具有有限长度。设 F 是 M 上的 q-良过滤。存在 \(n_{1} \in Z\)，使 \(F_{n_{1}} = M\)，从而 \(F_{n} = M\) 对 \(n \leqslant n_{1}\) 成立；因此我们通过置下式定义元素 \(H_{M,F}\) 属于 \(\mathbf{Z}((T))\)：

\[
\mathrm{H} _ {\mathrm{M}, \mathrm{F}} = \sum_ {n \in \mathbf {Z}} \operatorname{long} _ {\mathrm{A} / \mathfrak {q}} (\mathrm{F} _ {n} / \mathrm{F} _ {n + 1}). \mathrm{T} ^ {n} \in \mathbf {Z} ((\mathrm{T})) .\tag{8}
\]

定义 1. --- 我们称 H \(_{M,F}\) 为 A-模 M 的希尔伯特–塞缪尔级数（相对于 q-良过滤 F）。

映射 \(n\mapsto\text{long}_{\mathbf{A}}(\mathbf{F}_{n}/\mathbf{F}_{n+1})\) 通常称为 M 的希尔伯特–塞缪尔函数（相对于 F）。

这尤其适用于 q-进过滤的情形（\(F_{n} = q^{n}M\)）；此时置 \(H_{M,F} = H_{M,q}\)。因此有

\[
\mathrm{H} _ {\mathrm{M}, \mathfrak {q}} = \sum_ {n \in \mathbf {Z}} \operatorname{long} _ {\mathrm{A} / \mathfrak {q}} (\mathfrak {q} ^ {n} \mathrm{M} / \mathfrak {q} ^ {n + 1} \mathrm{M}). \mathrm{T} ^ {n}.\tag{9}
\]

命题 3.--- a) 若 F 是 M 上的 q-良过滤，则有

\[
\mathrm{H} _ {\mathrm{M}, \mathrm{F}} ^ {(1)} = \sum_ {n \in \mathbf {Z}} \operatorname{long} _ {\mathrm{A}} (\mathrm{M} / \mathrm{F} _ {n + 1}). \mathrm{T} ^ {n}.\tag{10}
\]

\begin{enumerate}
\def\labelenumi{\alph{enumi})}
\setcounter{enumi}{1}
\tightlist
\item
  若 F 和 F' 是 M 上的两个 q-良过滤，则存在整数 m，使 \(\mathrm{H}_{\mathbf{M},\mathrm{F}'}^{(1)} \geqslant \mathrm{T}^{m}\mathrm{H}_{\mathbf{M},\mathbf{F}}^{(1)}\)。
\end{enumerate}

部分 \(a)\) 立即由 \(\mathbf{H}_{\mathbf{M},\mathbf{F}}^{(1)}\) 的定义得到；部分 \(b)\) 由 \(a)\) 和引理 3 得到。

定理 2. --- 设 A 为诺特环，q 为 A 的理想，M 为有限生成 A-模，使 M/qM 非零且具有有限长度，F 为 M 上的 q-良过滤。

\begin{enumerate}
\def\labelenumi{\alph{enumi})}
\item
  存在唯一确定的整数 \(d \geqslant 0\) 和元素 \(\mathbf{R}\)，属于 \(\mathbf{Z}[\mathrm{T}, \mathrm{T}^{-1}]\)，使 \(\mathbf{R}(1) > 0\) 且 \(\mathrm{H}_{\mathrm{M,F}} = (1 - \mathrm{T})^{-d}\mathbf{R}\)。
\item
  整数 d 和 R(1) 与所选的 q-良过滤 F 无关。
\item
  考虑分次环 \(\operatorname{gr}(A)\)，其中 \(\operatorname{gr}_{n}(A) = \mathfrak{q}^{n}/\mathfrak{q}^{n+1}\)，以及分次 \(\operatorname{gr}(A)\)-模 \(\operatorname{gr}(M)\)，其中 \(\operatorname{gr}_{n}(M) = F_{n}/F_{n+1}\)。由于有 \(F_{n_{1}} = M\)，且 \(\mathfrak{q}F_{n} = F_{n+1}\) 对 \(n \geqslant n_{0}\) 成立，\(\operatorname{gr}(M)\) 由 \(\bigoplus_{n_{1} \leqslant n \leqslant n_{0}} \operatorname{gr}_{n}(M)\) 生成，因而是有限型的。此外，若 \((x_{1},\ldots,x_{r})\) 是 A-模 \(\mathfrak{q}\) 的有限生成系，则 \(\operatorname{gr}(A)\) 由 \(\operatorname{gr}_{0}(A)\) 以及 \(x_{i}\) 模 \(\mathfrak{q}^{2}\) 的类生成，因而同构于 \(H = (A/\mathfrak{q})[X_{1},\ldots,X_{r}]\) 的一个分次商环。由 no. 2 的 th. 1，有
\end{enumerate}

\[
(1 - \mathrm{T}) ^ {r} \mathrm{H} _ {\mathrm{M}, \mathrm{F}} \in \mathbf {Z} [ \mathrm{T}, \mathrm{T} ^ {- 1} ].
\]

有 \(H_{M,F} \neq 0\)，因此存在唯一确定的 \(d \in N\) 和 \(R \in Z[T, T^{-1}]\)，使 \(\mathrm{R}(1) > 0\) 且 \(\mathrm{H}_{\mathrm{M},\mathrm{F}} = (1 - \mathrm{T})^{-d}. R\)（no. 1 的引理 2）。

\begin{enumerate}
\def\labelenumi{\alph{enumi})}
\setcounter{enumi}{1}
\tightlist
\item
  令 F' 为另一个 q-良过滤，并类似地写成
\end{enumerate}

\[
\mathrm{H} _ {\mathbf {M}, \mathrm{F} ^ {\prime}} = (1 - \mathrm{T}) ^ {- d ^ {\prime}} \mathbf {R} ^ {\prime}.
\]

由 Prop. 3, b)，存在整数 m，使 \((1 - T)^{-d'-1}R' \geqslant T^{m}(1 - T)^{-d-1}R\)。根据 no. 1 的引理 2, b)，这蕴含 \(d' \geqslant d\)，并且若 \(d' = d\)，则 \(R'(1) \geqslant R(1)\)。交换 F 和 \(F'\) 的角色，得到 \(d = d'\) 和 \(R(1) = R'(1)\)。

备注 1. --- 用 \(a\) ) 的记号，写成 \(\mathbf{R} = \sum_{i \in \mathbf{Z}} a_i \mathbf{T}^i\)，并假设 \(d > 0\)。由 no. 1，关系 \(\mathrm{H}_{\mathrm{M},\mathrm{F}} = (1 - \mathrm{T})^{-d}\mathrm{R}\) 也可写成

\[
\operatorname{long} _ {\mathbf {A}} \left(\mathrm{F} _ {n} / \mathrm{F} _ {n + 1}\right) = \sum_ {i \in \mathbf {Z}} a _ {i} \left[ \begin{array}{c} n - i + d - 1 \\ d - 1 \end{array} \right] = \sum_ {i \leqslant n} a _ {i} \binom {n - i + d - 1} {d - 1}.\tag{11}
\]

类似地，由于 \(\mathbf{H}_{\mathbf{M},\mathbf{F}}^{(1)} = (1 - \mathbf{T})^{-d - 1}\mathbf{R}\)，有

\[
\operatorname{long} _ {\mathbf {A}} \left(\mathrm{M} / \mathrm{F} _ {n + 1}\right) = \sum_ {i \in \mathbf {Z}} a _ {i} \left[ \begin{array}{c} n - i + d \\ d \end{array} \right] = \sum_ {i \leqslant n} a _ {i} \binom {n - i + d} {d}.\tag{12}
\]

设 A 为诺特环，q 为 A 的理想，M 为有限生成 A-模，且 M/qM 具有有限长度。若 M ≠ qM，则根据定理 2, b)，存在整数 \(d_{q}(M) \geqslant 0\) 和 \(e_{q}(M) > 0\)，使得对于 M 上的每个 q-良过滤，都存在 R ∈ Z{[}T, T \(^{-1}\) {]}，满足

\[
\mathrm{H} _ {\mathbf {M}, \mathrm{F}} = (1 - \mathrm{T}) ^ {- d _ {\mathfrak {q}} (\mathbf {M})} \mathrm{R}, \quad \mathrm{R} (1) = e _ {\mathfrak {q}} (\mathbf {M}).
\]

若 M = qM，则按约定置 \(d_{q}(M) = -\infty\)，\(e_{q}(M) = 0\)。

备注 2. --- 说 M/qM 具有有限长度，意味着

\[
\operatorname{Supp} (\mathrm{M} / \mathfrak {q} \mathrm{M}) = \operatorname{Supp} (\mathrm{M}) \cap \mathrm{V} (\mathfrak {q})
\]

由极大理想组成（IV, § 2, no. 5, prop. 7）。我们将在下文（no. 4, corollary to th. 3）看到，\(d_{\mathfrak{q}}(\mathbf{M})\) 是数 \(\dim_{\mathbf{A}_{\mathfrak{m}}}(M_{\mathfrak{m}})\) 的最小上界，其中 m 遍历集合 \(\operatorname{Supp}(\mathbf{M}) \cap \mathbf{V}(\mathfrak{q})\)。

推论。— 设 A 为诺特环，q 为 A 的理想，M 为有限生成 A-模，使 M/qM 具有有限长度，F 为 M 上的 q-良过滤。

\begin{enumerate}
\def\labelenumi{\alph{enumi})}
\tightlist
\item
  要有 \(d_{\mathfrak{q}}(\mathbf{M}) \leqslant 0\)，必要且充分的条件是序列 \((\mathfrak{q}^{n}\mathbf{M})\) 稳定，或等价地，序列 \((\mathrm{F}_{n})\) 稳定。于是对于所有充分大的 n，
\end{enumerate}

\[
\operatorname{long} \left(\mathrm{M} / \mathrm{F} _ {n + 1}\right) = \operatorname{long} \left(\mathrm{M} / \mathfrak {q} ^ {n + 1} \mathrm{M}\right) = e _ {\mathfrak {q}} (\mathrm{M}).
\]

\begin{enumerate}
\def\labelenumi{\alph{enumi})}
\setcounter{enumi}{1}
\tightlist
\item
  假设 \(d_{\mathfrak{q}}(\mathbf{M}) > 0\)。则有
\end{enumerate}

\begin{enumerate}
\def\labelenumi{(\arabic{enumi})}
\setcounter{enumi}{12}
\tightlist
\item
\end{enumerate}

\[
\operatorname{long} _ {\mathbf {A}} \left(\mathrm{F} _ {n} / \mathrm{F} _ {n + 1}\right) = e _ {\mathfrak {q}} (\mathrm{M}) n ^ {d _ {\mathfrak {q}} (\mathrm{M}) - 1} / \left(d _ {\mathfrak {q}} (\mathrm{M}) - 1\right)! + \rho_ {n} n ^ {d _ {\mathfrak {q}} (\mathrm{M}) - 2},\tag{14}
\]

\[
\operatorname{long} _ {\mathbf {A}} (\mathrm{M} / \mathrm{F} _ {n + 1}) = e _ {\mathfrak {q}} (\mathbf {M}) n ^ {d _ {\mathfrak {q}} (\mathbf {M})} / d _ {\mathfrak {q}} (\mathbf {M})! + \sigma_ {n} n ^ {d _ {\mathfrak {q}} (\mathbf {M}) - 1},
\]

其中 \(\rho_{n}\) 和 \(\sigma_{n}\) 在 \(n\) 无限增大时趋于一个极限。

这立即由 th. 2 和 no. 1 的公式 (2) 得到。

备注 3. --- 假设 q 包含于 A 的根基中。于是根据 Nakayama 引理，序列 (q \(^{n}\) M) 稳定，当且仅当对于充分大的 n 有 q \(^{n}\) M = 0。由推论的 a) 部分随即得到，d \(_{q}\) (M) ≤ 0 当且仅当 M 具有有限长度，并且此时 e \(_{q}\) (M) = long \(_{A}\) (M)。

命题 4. — 设 A 为诺特环，\(x_{1},...,x_{r}\) 是 A 中的元素，\(\mathfrak{x}\) 是它们生成的理想，M 是有限生成 A-模，且 \(M/\mathfrak{x}M\) 非零并具有有限长度。

\begin{enumerate}
\def\labelenumi{\alph{enumi})}
\item
  有 \(d_{\mathfrak{x}}(\mathbf{M}) \leqslant r\)。
\item
  若 \(d_{\mathfrak{x}}(M) = r\)，则 \(e_{\mathfrak{x}}(M) \leqslant \text{long}_{A}(M/\mathfrak{x}M)\)。
\item
  若序列 \((x_{1},...,x_{r})\) 是 M 的完全截切序列（A, X, p. 157），则 \(d_{\mathfrak{x}}(\mathbf{M}) = r\)，且 \(e_{\mathfrak{x}}(\mathbf{M}) = \text{long}_{\mathbf{A}}(\mathbf{M}/\mathfrak{x}\mathbf{M})\)。若 \(x_{i}\) 属于 A 的根基，则逆命题成立。
\end{enumerate}

令 H 为多项式环 \((\mathbf{A} / \mathfrak{x})[\mathbf{X}_1,\dots ,\mathbf{X}_r]\)。我们赋予 \(\mathbf{G} = \bigoplus_{n}\mathfrak{x}^{n}\mathbf{M} / \mathfrak{x}^{n + 1}\mathbf{M}\)

一个分次 H-模结构，其中 \((\mathrm{X}_{i})_{\mathrm{G}}\) 是乘以 \(x_{i}\) 模 \(\mathfrak{x}^{2}\) 的类的映射。采用 No. 2 中的记号 \(P_{G}, Q_{G}, c_{G}\)，有 \(H_{M,\mathfrak{x}} = P_{G}\)，因而 \((1 - T)^{-d_{\mathfrak{x}}(M)}R = (1 - T)^{-r}Q_{G}\)，其中 \(R(1) = e_{\mathfrak{x}}(M) > 0\)，且 \(Q_{G}(1) = c_{G}\)。

因此，要么 \(d_{\mathfrak{x}}(\mathbf{M}) < r\) 且 \(c_{\mathbf{G}} = 0\)，要么 \(d_{\mathfrak{x}}(\mathbf{M}) = r\) 且 \(c_{\mathbf{G}} = e_{\mathfrak{x}}(\mathbf{M})\)。此外，由 no. 2 的 prop. 1，有 \(c_{\mathbf{G}} \leqslant \text{long}(\mathbf{M}/\mathfrak{x}\mathbf{M})\)，且等号成立当且仅当典范同态 \(\mathrm{A}/\mathfrak{x}[\mathrm{X}_{1}, ..., \mathrm{X}_{r}] \otimes_{\mathrm{A}/\mathfrak{x}} \mathrm{M}/\mathfrak{x}\mathrm{M} \to \bigoplus_{n}\mathfrak{x}^{n}\mathrm{M}/\mathfrak{x}^{n+1}\mathrm{M}\) 是双射。

根据 A, X, p. 160 的 th. 1，这就证明了命题。

命题 5. --- 设 \(0 \to \mathbf{M}' \to \mathbf{M} \to \mathbf{M}'' \to 0\) 是诺特环 A 上有限生成模的正合序列，q 是 A 的理想。

\begin{enumerate}
\def\labelenumi{\alph{enumi})}
\item
  若 \(M/\mathfrak{q}M\) 具有有限长度，其充要条件是 \(M'/\mathfrak{q}M'\) 和 \(M''/\mathfrak{q}M''\) 也具有有限长度。
\item
  假设 M/qM 具有有限长度。则属于以下三种情形之一：
\end{enumerate}

\[
1) d _ {\mathfrak {q}} (\mathbf {M}) = d _ {\mathfrak {q}} \left(\mathbf {M} ^ {\prime}\right) > d _ {\mathfrak {q}} \left(\mathbf {M} ^ {\prime \prime}\right) and e _ {\mathfrak {q}} (\mathbf {M}) = e _ {\mathfrak {q}} \left(\mathbf {M} ^ {\prime}\right),
\]

\begin{enumerate}
\def\labelenumi{\arabic{enumi})}
\setcounter{enumi}{1}
\tightlist
\item
  \(d_{\mathfrak{q}}(\mathbf{M}) = d_{\mathfrak{q}}(\mathbf{M}^{\prime \prime}) > d_{\mathfrak{q}}(\mathbf{M}^{\prime})\) 且 \(e_{\mathfrak{q}}(\mathbf{M}) = e_{\mathfrak{q}}(\mathbf{M}^{\prime \prime}),\)
\end{enumerate}

\[
3) d _ {\mathfrak {q}} (\mathbf {M}) = d _ {\mathfrak {q}} \left(\mathbf {M} ^ {\prime}\right) = d _ {\mathfrak {q}} \left(\mathbf {M} ^ {\prime \prime}\right) and e _ {\mathfrak {q}} (\mathbf {M}) = e _ {\mathfrak {q}} \left(\mathbf {M} ^ {\prime}\right) + e _ {\mathfrak {q}} \left(\mathbf {M} ^ {\prime \prime}\right).
\]

\begin{enumerate}
\def\labelenumi{\alph{enumi})}
\item
  有 \(\operatorname{Supp}(\mathbf{M}) = \operatorname{Supp}(\mathbf{M}') \cup \operatorname{Supp}(\mathbf{M}'')\)，断言由备注 2 得到。
\item
  赋予 M 一个 q-良过滤 F（例如 q-进过滤），赋予 M'' 商过滤 F''，赋予 M' 诱导过滤 F'。过滤 F' 和 F'' 都是 q-良过滤（III, § 3, no. 1, prop. 1）。于是对于每个 n，有 A-模的正合序列
\end{enumerate}

\[
0 \rightarrow \mathrm{F} _ {n} ^ {\prime} / \mathrm{F} _ {n + 1} ^ {\prime} \rightarrow \mathrm{F} _ {n} / \mathrm{F} _ {n + 1} \rightarrow \mathrm{F} _ {n} ^ {\prime \prime} / \mathrm{F} _ {n + 1} ^ {\prime \prime} \rightarrow 0
\]

（III, § 2, no. 4, prop. 2），因而有 \(\mathrm{H}_{\mathbf{M},\mathbf{F}} = \mathrm{H}_{\mathbf{M}',\mathbf{F}'} + \mathrm{H}_{\mathbf{M}'',\mathbf{F}''}\)，或等价地

\[
(1 - \mathrm{T}) ^ {- d _ {\mathfrak {q}} (\mathrm{M})} \mathrm{R} = (1 - \mathrm{T}) ^ {- d _ {\mathfrak {q}} (\mathrm{M} ^ {\prime})} \mathrm{R} ^ {\prime} + (1 - \mathrm{T}) ^ {- d _ {\mathfrak {q}} (\mathrm{M} ^ {\prime \prime})} \mathrm{R} ^ {\prime \prime},
\]

其中 \(R,R',R''\in\mathbf{Z}[T,T^{-1}]\)，\(R(1)=e_{\mathfrak{q}}(\mathbf{M})\)，\(R'(1)=e_{\mathfrak{q}}(\mathbf{M}')\)，\(R''(1)=e_{\mathfrak{q}}(\mathbf{M}'')\)。断言 b) 立即由此得到。

